\section{分岐曲線の既約性と次数二十四}
\label{sec:septic-branch-integrality-v30}

有限性の証明に使った特殊曲面には三本の収縮二次曲線が残る。
その退化を用いると、一般の分岐曲線の既約性も決定できる。
要点は、特殊な分岐源の中で $L$ と正に交わる成分が一本しかないことである。

\begin{theorem}[一般の分岐曲線]
\label{thm:septic-branch-integrality-v30}
同じ選択成分 $\mathcal Y_0$ のある非空開部分で、次が成り立つ。
\begin{enumerate}
\item 標準七次射 $f:X\to\mathbb P^2$ は有限平坦であり、分岐源 $R_f$ は幾何学的に整である。
\item 被約分岐曲線 $B_f\subset\mathbb P^2$ は幾何学的に整な二十四次曲線である。
写像 $R_f\to B_f$ は双有理であり、一般分岐点での慣性は互換である。
\item $p_a(R_f)=37$ であり、
\[
 \delta(B_f)=216+\delta(R_f)
\]
が成り立つ。ただし $\delta$ は正規化による算術種数の減少を表す。
次節の定理\ref{thm:septic-ramification-smoothness-v30}では、さらに
$R_f$ が滑らかであることを証明し、正規化種数三十七と $\delta(B_f)=216$ を得る。
\end{enumerate}
\end{theorem}

\subsection{既約な残余六次曲線を持つ特殊点}

六変数部分族で別の有理点
\[
 (a,c,h,j,n,q)=(1,2,3,4,5,6),\qquad s=1
\]
を選び、これを $A_\#$ と書く。この点は $A_0$ を通る既約な六変数族にあるため、
前と同じ選択成分 $\mathcal Y_0$ に属する。
ここでは
\[
 M_\#=\begin{pmatrix}-1&1&2\\3&-1&4\\5&6&-1\end{pmatrix},\qquad
 \det M_\#=92.
\]
$F=M_\#x^{[2]}$ とし、残余critical sexticを
\[
 J(x)=\det\bigl(\operatorname{diag}(F)-2\operatorname{diag}(x)M_\#\operatorname{diag}(x)\bigr)
\]
と置く。明示式は
\begin{align*}
 J={}&15x_0^6-22x_0^4x_1^2-67x_0^4x_2^2-43x_0^2x_1^4
      -502x_0^2x_1^2x_2^2-110x_0^2x_2^4\\
     &+6x_1^6-59x_1^4x_2^2-138x_1^2x_2^4+8x_2^6.
\end{align*}

\begin{lemma}[分岐源の被約な特殊ファイバー]
\label{lem:septic-special-ramification-v30}
$A_\#$ で $C=V(\det Q)$ と $V(J)$ はともに滑らかである。
$V(J)$ の狭義変換を $R_0$ とすると
\[
 R_{f_\#}=R_0+T_0+T_1+T_2
\]
はすべての成分の係数が一の被約除数である。
$R_0$ は滑らかな種数十の曲線であり、
\[
 [R_0]=6H-\sum_{p\in D_{\rm coord}}E_p,\qquad
 L\cdot R_0=24,\qquad L\cdot T_i=0,
\]
ただし $D_{\rm coord}=\bigcup_iV(x_i,F_i)$ は六点からなる。
\end{lemma}
\begin{proof}
三つの二次曲線と十八基点の横断性は、$M_\#$ の全成分と二次小行列式の非消滅から従う。
$C$ と $J$ の滑らかさは $\mathbb F_{101}$ 上で証明される。
両曲線について三つのアフィンチャートすべてで、三つの斉次偏微分からなるイデアルの
Gr\"obner基底が $\{1\}$ となる。
これは射影的な特異点がないことを示し、標数零での滑らかさも従う。
計算は \nolinkurl{septic_ramification_v30.py} と随伴JSONに記録する。

既存の恒等式
\[
 \det(dg)=-5F_0F_1F_2J
\]
を使う。$J$ は十二個の $F_i=F_j=0$ を避け、
六個の $x_i=F_i=0$ を通る。
前者は $F_i,F_j$ による交点の二乗座標を代入して非零を確かめ、
後者は $J|_{x_i=0}$ が $F_i|_{x_i=0}$ で割り切れることから従う。
滑らかさにより後者の六点での重複度は一である。
従って、十二点では $F_0F_1F_2J$ の重複度は二、残る六点でも二である。

単純基点の一度のblow-upで、平面上のJacobian除数の狭義変換に対する例外係数は
基点重複度から二を引いた値となる。
これは $g_i=uh_i$ と書くblow-upチャートで
$dg_0\wedge dg_1\wedge dg_2$ を計算することによっても得られる。
ここでは例外成分は残らず、主張した四成分の和がちょうど $R_{f_\#}$ である。
$J$ は滑らかな平面六次曲線なので、その狭義変換も滑らかで種数は十である。
交点数は
\[
 L\cdot R_0=5\cdot6-6=24
\]
である。
\end{proof}

\begin{proof}[定理\ref{thm:septic-branch-integrality-v30}の証明]
十八基点が被約である良い開部分上で相対blow-upを行い、
滑らかな曲面族 $\mathcal X$ と射 $\mathcal f$ を考える。
相対微分の行列式は各ファイバーで零でない切断なので、
分岐源 $\mathcal R$ は相対有効Cartier除数であり、底上平坦である。

幾何学的一般ファイバーにおいて、分岐源が複数成分を持つか、
ある成分の係数が二以上であると仮定する。
必要なら底の有限拡大を行って成分を定義し、
$A_\#$ への離散付値環を通じてそれらの有効サイクルを特殊化する。
特殊化は有効性、$L$ との交点数、和としての $\mathcal R$ を保つ。
一般ファイバーでは定理\ref{thm:septic-finiteness-v30}により $L$ は豊富なので、
各成分は $L$ と正に交わる。

一方、補題\ref{lem:septic-special-ramification-v30}の特殊ファイバーでは、
$L$ と正に交わる成分は $R_0$ のみである。
従って各一般成分の特殊化は $R_0$ を含まなければならない。
しかし特殊ファイバーにおける $R_0$ の係数は一である。
このため一般成分は一本、その係数も一に限られる。
幾何学的一般分岐源は整であり、平坦な固有族における幾何学的整性の開性により、
底を縮めれば全ファイバーでこの性質が成り立つ。

前節で証明した単純foldの持続開部分と交わらせる。
そのfoldの標的近傍で、有限平坦代数は分岐する階数二の因子と
\'etaleな五つの階数一因子に分かれる。
従ってその近傍の分岐曲線は滑らかであり、分岐源からはただ一つの点が一対一に写る。
既約な $R_f$ からその像 $B_f$ への有限写像の一般次数は、これにより一である。
$R_f$ は係数一の分岐源なので、一般分岐指数は二である。
既存の押出し計算
\[
 \deg f_*R_f=L\cdot R_f=24
\]
は、ここで被約な既約分岐曲線そのものの次数二十四を与える。

最後に $R_f\sim12H-2\sum E_i$ から
\[
 R_f^2=72,\qquad K_XR_f=0,\qquad p_a(R_f)=1+72/2=37.
\]
$B_f$ は二十四次なので $p_a(B_f)=23\cdot22/2=253$ である。
両曲線は双有理で同じ正規化を持つため
\[
 \delta(B_f)-\delta(R_f)=253-37=216
\]
となる。
\end{proof}

\begin{corollary}[六変数収縮部分族の分岐]
\label{cor:septic-contracted-branch-v30}
六変数部分族の一般点では、残余分岐源は滑らかな種数十の曲線である。
その像は既約な二十四次分岐曲線であり、残余分岐源からの写像はその正規化である。
\end{corollary}
\begin{proof}
補題\ref{lem:septic-special-ramification-v30}により、
残余六次曲線が滑らかである条件はこの既約な六変数族の非空開条件である。
前節の単純foldの持続論法は、滑らかな六変数部分族そのものを底にしても適用できる。
従って、同じ族の別の非空開部分で単純foldと五つの\'etale点からなる完全なファイバーが存在する。
二つの開部分は交わる。
唯一の正次数の分岐源は残余六次曲線なので、同じ次数一の議論により結論を得る。
\end{proof}

\paragraph{次節での強化.}
定理\ref{thm:septic-ramification-smoothness-v30}では、同じ選択成分内で
$R_f$ が滑らかであることを証明する。
これにより $R_f$ 自身が種数三十七の正規化となり、$\delta(B_f)=216$ が確定する。
さらに正規化写像の微分零因子の長さと双対曲線の次数を求める。
通常の節点・尖点だけが生じるという特異点型の限定は、これらとは別に残る問題である。
